\begin{center}
    {\LARGE \textbf{2008物理化学}} \\[2ex]
    {\large 时量180分钟 \quad 满分150分}
\end{center}

\vspace{0.5cm}
\noindent\textbf{注意事项:}
\begin{enumerate}[label=\arabic*., parsep=0pt, itemsep=0pt]
    \item 答题前,考生先将自己的姓名、学号、考场号、座位号填写在答题卡上,将条形码准确粘贴在考生信息条形码粘贴区。
    \item 考试结束后,将本试卷与答题纸一并收回。
    \item 回答各个问题时,将答案写在答题纸上,写在本试卷上无效。
\end{enumerate}

\vspace{0.5cm}
\noindent\textbf{1. (15 pts.)}

$1\ \mathrm{mol}$单原子分子理想气体经过一个绝热不可逆过程到达终态,该终态的温度为$273\ \mathrm{K}$,压力为$p^\ominus$,熵值为$S_\mathrm{m}^\ominus(273\ \mathrm{K})=188.3\ \mathrm{J\cdot K^{-1}\cdot mol^{-1}}$。

已知该过程的$\Delta S_\mathrm{m}=20.92\ \mathrm{J\cdot K^{-1}\cdot mol^{-1}}$,$W=-1255\ \mathrm{J}$。
\begin{enumerate}[label=(\arabic*), itemsep=1ex]
    \item 求始态的$p_1,V_1,T_1$；
    \item 求气体的$\Delta U,\Delta H,\Delta G$。
\end{enumerate}

\vspace{0.5cm}
\noindent\textbf{2. (15 pts.)}

$300\ \mathrm{K}$时,试计算(1) 从大量的等物质的量的$\ce{C_{2}H_{4}Br_{2}}$和$\ce{C_{3}H_{6}Br_{2}}$的理想混合物中分离出$1\ \mathrm{mol}$纯$\ce{C_{2}H_{4}Br_{2}}$过程的$\Delta G_1$,(2) 若混合物中各含$2\ \mathrm{mol}$ $\ce{C_{2}H_{4}Br_{2}}$和$\ce{C_{3}H_{6}Br_{2}}$,从中分离出$1\ \mathrm{mol}$纯$\ce{C_{2}H_{4}Br_{2}}$时,$\Delta G_2$又为多少。

\vspace{0.5cm}
\noindent\textbf{3. (15 pts.)}

液态氨和固态氨的蒸气压与温度的关系分别为
\[
\ln[p(\mathrm{l})/\mathrm{Pa}]=24.38-\frac{3063\ \mathrm{K}}{T}
\]
\[
\ln[p(\mathrm{s})/\mathrm{Pa}]=27.92-\frac{3754\ \mathrm{K}}{T}
\]
试求(1) 三相点的温度与压力；
(2) 三相点的蒸发热、升华热和熔化热。

\vspace{0.5cm}
\noindent\textbf{4. (15 pts.)}

实验证明:两块没有氧化膜的光滑洁净的金属表面紧靠在一起时,它们会自动地黏合在一起。假定外层空间的大气压力约为$1.013\times10^{-9}\ \mathrm{Pa}$,温度的影响暂不考虑,当两个镀铬的宇宙飞船由地面进入外层空间对接时,它们能否自动地黏合在一起。

已知$\Delta_\mathrm{f}G_\mathrm{m}^\ominus(\ce{Cr_{2}O_{3}(s)},298.15\ \mathrm{K})=-1079\ \mathrm{kJ\cdot mol^{-1}}$。设外层空间的温度为$298\ \mathrm{K}$,空气的组成与地面相同($\ce{O_{2}}$占五分之一)。

\vspace{0.5cm}
\noindent\textbf{5. (15 pts.)}

在烧瓶中装有$\ce{N_{2}O_{5}}$和$\ce{CCl_{4}}$的溶液,把该烧瓶放入$40^\circ\mathrm{C}$恒温槽中时,$\ce{N_{2}O_{5}}$开始按一级反应分解成$\ce{NO_{2}}$、$\ce{N_{2}O_{4}}$和$\ce{O_{2}}$,其中$\ce{NO_{2}}$、$\ce{N_{2}O_{4}}$留在溶液中,而$\ce{O_{2}}$则通过一根狭管在$40^\circ\mathrm{C}$和某一定压力下收集在$\ce{CCl_{4}}$的液面上。当收集$10.75\ \mathrm{mL}$气体时开始进行测定($t=0$),在$t=2400\ \mathrm{s}$时,测得气体体积为$29.65\ \mathrm{mL}$,经很长时间后($t=\infty$)体积为$45.50\ \mathrm{mL}$。试求
\begin{enumerate}[label=(\arabic*), itemsep=1ex]
    \item 速率常数；
    \item $40^\circ\mathrm{C}$时反应的半衰期；
    \item $4800\ \mathrm{s}$后收集到多少体积的氧。
\end{enumerate}

\vspace{0.5cm}
\noindent\textbf{6. (15 pts.)}

电池 $\ce{Zn|ZnCl_{2}(0.05\ mol\cdot kg^{-1})}|\ce{AgCl(s)}|\ce{Ag(s)}$的电动势与温度的关系为
\[
E/V=1.015-4.92\times10^{-4}(T/K-298)
\]
\begin{enumerate}[label=(\arabic*), itemsep=1ex]
    \item 写出电池反应并计算该电池在$298\ \mathrm{K}$时的电动势。
    \item 计算$298\ \mathrm{K}$时上述电池可逆输出$1\ \mathrm{F}$电荷量时,电池反应的$\Delta_\mathrm{r}G_\mathrm{m}$、$\Delta_\mathrm{r}S_\mathrm{m}$、$\Delta_\mathrm{r}H_\mathrm{m}$、$\Delta_\mathrm{r}U_\mathrm{m}$、$\Delta_\mathrm{r}F_\mathrm{m}$以及该电池的热效应。
    \item 若将该电池短路,上述各值为多少?
    \item 若该电池以$0.8\ \mathrm{V}$工作电压放电,上述各值又为多少?
\end{enumerate}